%=============================================================================
% Temporal_Sector_Analysis.tex
% Critical Analysis: The Dual Character of the DFD psi-Field
% Constraint-Like (Elliptic) vs. Propagating (Hyperbolic) Behavior
% Date: 2026-03-23
%=============================================================================
\documentclass[12pt,a4paper]{article}
\usepackage[utf8]{inputenc}
\usepackage{amsmath,amssymb,amsfonts,amsthm}
\usepackage{hyperref}
\usepackage{booktabs}
\usepackage{geometry}
\usepackage{xcolor}
\usepackage{tcolorbox}
\usepackage{array}
\geometry{margin=2.5cm}

\newtheorem{theorem}{Theorem}[section]
\newtheorem{proposition}[theorem]{Proposition}
\newtheorem{lemma}[theorem]{Lemma}
\newtheorem{corollary}[theorem]{Corollary}
\newtheorem{remark}{Remark}[section]

\definecolor{dfdgreen}{RGB}{0,128,0}
\definecolor{dfdblue}{RGB}{0,50,180}
\definecolor{dfdred}{RGB}{200,0,0}

\newcommand{\psifield}{\psi}
\newcommand{\order}[1]{\mathcal{O}(#1)}

\title{The Dual Character of the DFD $\psi$-Field:\\
Temporal Sector Analysis and the Validity\\
of the Quasi-Static Constraint Approximation\\[6pt]
\large Correcting and Refining the i32--i34 Characterization\\[6pt]
\normalsize DFD Research Programme}
\author{Cross-Reference Analysis}
\date{23 March 2026}

\begin{document}
\maketitle
\tableofcontents
\newpage

%=============================================================================
\section{The Issue: $\psi$ Is Not Literally $A_0$}
\label{sec:issue}
%=============================================================================

The QG iteration analyses (i32--i34) established a powerful analogy between
the DFD $\psi$-field and the Coulomb potential $A_0$ in QED. The key claims
were:
\begin{enumerate}
\item $\psi$ has no free propagator (i32, Theorem 1).
\item $\psi$ is a ``constraint field'' determined by matter sources.
\item Quantization proceeds via branch-dependent classical solutions, not
      Fock space.
\end{enumerate}

These claims are \textbf{correct in the quasi-static limit}, but they require
careful qualification because the full DFD field equation is:
\begin{equation}\label{eq:full-dfd}
\boxed{\frac{1}{c^2}\frac{\partial^2\psi}{\partial t^2}
- \nabla\cdot\!\left[\mu\!\left(\frac{|\nabla\psi|^2}{a_\star^2}\right)
\nabla\psi\right]
= -\frac{8\pi G}{c^2}(\rho - \bar\rho)}
\end{equation}

This equation has a \textbf{second time derivative}. It is a nonlinear wave
equation, not a constraint equation. In contrast, $A_0$ in Coulomb-gauge QED
satisfies Gauss's law $\nabla^2 A_0 = -\rho/\epsilon_0$, which is purely
elliptic---\textbf{no time derivatives whatsoever}. The Lagrangian for QED
in Coulomb gauge contains no $\dot{A}_0^2$ term; $A_0$ is truly
non-dynamical.

The purpose of this analysis is to establish:
\begin{enumerate}
\item Precisely when the quasi-static (constraint) approximation is valid.
\item The propagation properties of $\psi$ when it is \emph{not} quasi-static.
\item That for \textbf{all} laboratory quantum gravity experiments (Page-Geilker,
      BMV, etc.), the quasi-static limit is valid by many orders of magnitude.
\item The correct characterization of $\psi$: a field with \emph{both}
      constraint-like and propagating behavior, with the former dominating in
      all experimentally relevant quantum regimes.
\end{enumerate}


%=============================================================================
\section{The Full Equation: Classification and Characteristics}
\label{sec:classification}
%=============================================================================

\subsection{PDE type of the full equation}

Equation~(\ref{eq:full-dfd}) is a second-order quasilinear PDE. The principal
part (highest-derivative terms) is:
\begin{equation}
\frac{1}{c^2}\partial_t^2\psi
- \mu(\chi)\nabla^2\psi
- \mu'(\chi)\frac{2}{a_\star^2}(\partial_i\psi)(\partial_j\psi)\partial_i\partial_j\psi = \ldots
\end{equation}
where $\chi = |\nabla\psi|^2/a_\star^2$.

The principal symbol is:
\begin{equation}
P(\xi_0, \boldsymbol{\xi}) = \frac{\xi_0^2}{c^2}
- \mu(\chi)|\boldsymbol{\xi}|^2
- \frac{2\chi\mu'(\chi)}{|\nabla\psi|^2}(\nabla\psi\cdot\boldsymbol{\xi})^2
\end{equation}

The characteristics of this equation are the surfaces $P(\xi_0,\boldsymbol{\xi})=0$.
The equation is \textbf{hyperbolic} when the spatial part of the principal
symbol is positive-definite, which requires $\mu > 0$ and
$\mu + 2\chi\mu' > 0$. For DFD's $\mu(\chi) = \chi/(1+\chi)$:
\begin{align}
\mu &= \frac{\chi}{1+\chi} > 0 \quad\text{for }\chi > 0, \\
\mu + 2\chi\mu' &= \frac{\chi}{1+\chi} + \frac{2\chi}{(1+\chi)^2}
= \frac{\chi(3+\chi)}{(1+\chi)^2} > 0 \quad\text{for }\chi > 0.
\end{align}

Therefore: \textbf{on any non-trivial background ($\chi > 0$), the full DFD
equation is hyperbolic}. It is a wave equation with finite propagation speed.

At $\chi = 0$ (zero background gradient), the principal symbol degenerates:
$\mu(0) = 0$, and the spatial derivatives drop out entirely, leaving only
$\partial_t^2\psi/c^2 = \text{source}$. This is the cuscuton-like degeneracy
identified in i32.

\subsection{The sound speed}

On a non-trivial background $\bar\psi$ with $\bar\chi = |\nabla\bar\psi|^2/a_\star^2 > 0$, perturbations $\varphi = \psi - \bar\psi$ propagate with sound speed (i32, Theorem~2):
\begin{equation}\label{eq:cs}
\boxed{c_s^2 = \frac{1+\bar\chi}{3+\bar\chi}}
\end{equation}

The bounds are:
\begin{center}
\begin{tabular}{lcc}
\toprule
\textbf{Regime} & $\bar\chi$ & $c_s^2$ \\
\midrule
Deep MOND (weak field) & $\bar\chi \to 0$ & $1/3$ \\
Transition & $\bar\chi = 1$ & $1/2$ \\
Newtonian (strong field) & $\bar\chi \to \infty$ & $1$ \\
\bottomrule
\end{tabular}
\end{center}

In physical units:
\begin{equation}
\frac{c}{\sqrt{3}} \leq c_s \leq c
\end{equation}

The $\psi$-field propagates \textbf{subluminally everywhere}. In the Newtonian
regime ($\bar\chi \gg 1$), $c_s \to c$. In the deep-MOND regime
($\bar\chi \ll 1$), $c_s \to c/\sqrt{3} \approx 0.577\,c$.

\begin{remark}
The subluminality $c_s \leq c$ is crucial for causality. The fact that
$c_s < c$ for all finite $\bar\chi$ means $\psi$-perturbations travel
\emph{slower} than light. They are always ``behind'' photon signals. This
has profound implications for the constraint interpretation.
\end{remark}


%=============================================================================
\section{The Quasi-Static Limit: When $\psi$ Becomes a Constraint}
\label{sec:quasi-static}
%=============================================================================

\subsection{Formal derivation}

The quasi-static limit sets $\partial_t^2\psi/c^2 \approx 0$ in
Eq.~(\ref{eq:full-dfd}), yielding the elliptic equation:
\begin{equation}\label{eq:elliptic}
\nabla\cdot\!\left[\mu\!\left(\frac{|\nabla\psi|^2}{a_\star^2}\right)
\nabla\psi\right] = -\frac{8\pi G}{c^2}(\rho - \bar\rho)
\end{equation}

This is the standard AQUAL (AQUAdratic Lagrangian) equation of Bekenstein
\& Milgrom (1984). It is \textbf{elliptic}: given boundary conditions and
a source $\rho(\mathbf{x})$, there exists a unique solution
$\psi[\rho](\mathbf{x})$. The field is instantaneously determined by the
matter distribution, exactly like $A_0$ in Coulomb-gauge QED.

\subsection{Validity criterion}

The quasi-static approximation is valid when:
\begin{equation}\label{eq:qs-criterion}
\boxed{\frac{|\partial_t^2\psi/c^2|}{|\nabla\cdot[\mu\nabla\psi]|} \ll 1}
\end{equation}

For a system of characteristic size $L$ and characteristic timescale $T$
(the timescale on which the matter distribution changes):
\begin{equation}
\frac{\partial_t^2\psi}{c^2} \sim \frac{\psi}{c^2 T^2},
\qquad
\nabla\cdot[\mu\nabla\psi] \sim \frac{\mu\,\psi}{L^2}
\end{equation}

The ratio is:
\begin{equation}\label{eq:qs-ratio}
\frac{|\partial_t^2\psi/c^2|}{|\nabla\cdot[\mu\nabla\psi]|}
\sim \frac{L^2}{\mu\, c^2 T^2}
= \frac{1}{\mu}\left(\frac{L}{cT}\right)^2
\end{equation}

The quasi-static limit is valid when:
\begin{equation}
\frac{L}{cT} \ll \sqrt{\mu}
\end{equation}

Since $\mu \leq 1$, a sufficient condition is:
\begin{equation}\label{eq:qs-sufficient}
\boxed{L \ll c\,T}
\end{equation}

That is: \textbf{the quasi-static limit is valid whenever the system size
is much smaller than the distance light travels in one dynamical timescale}.
Equivalently: the $\psi$-propagation time across the system
($t_\psi \sim L/c_s$) is much shorter than the timescale on which the
source changes ($T$).

More precisely, using the sound speed:
\begin{equation}\label{eq:qs-precise}
\frac{t_\psi}{T} = \frac{L}{c_s\,T}
= \frac{L}{c\,T}\cdot\frac{c}{c_s}
= \frac{L}{c\,T}\cdot\sqrt{\frac{3+\bar\chi}{1+\bar\chi}} \ll 1
\end{equation}

Since $c/c_s \leq \sqrt{3}$, the quasi-static limit is valid whenever
$L \ll c\,T/\sqrt{3}$, i.e., whenever the system is non-relativistic.

\begin{theorem}[Quasi-Static Validity]
\label{thm:qs}
The quasi-static approximation for DFD's $\psi$-field is valid to accuracy
$\epsilon$ whenever:
\begin{equation}
\frac{L^2}{c_s^2\,T^2} < \epsilon
\end{equation}
where $L$ is the spatial extent of the mass distribution and $T$ is the
timescale on which the mass distribution changes. This is satisfied for
\textbf{all non-relativistic systems}, including all laboratory experiments,
solar system dynamics, and galactic dynamics.
\end{theorem}


%=============================================================================
\section{The Analogy with AQUAL/MOND: Standard Practice}
\label{sec:aqual}
%=============================================================================

The quasi-static limit of DFD is precisely the AQUAL equation
(\ref{eq:elliptic}), which is the standard starting point for
\textbf{all} MOND phenomenology:

\begin{itemize}
\item \textbf{Galaxy rotation curves}: Masses are non-relativistic ($v/c \sim 10^{-3}$). The dynamical timescale is the orbital period $T \sim 10^8$ years. The system size is $L \sim 30$ kpc. The ratio $L/(cT) \sim 10^{-3}$. Quasi-static is excellent.

\item \textbf{Galaxy clusters}: $v/c \sim 10^{-2}$. Still deeply non-relativistic. Quasi-static valid.

\item \textbf{Solar system}: $v/c \sim 10^{-4}$. Quasi-static is superb.

\item \textbf{Binary pulsars}: $v/c \sim 10^{-3}$ for the orbital velocity, but the gravitational wave emission requires the time-dependent equation. This is where the quasi-static limit begins to receive corrections.

\item \textbf{Compact object mergers}: $v/c \sim 0.3$--$0.5$. The quasi-static limit \textbf{fails}. The full hyperbolic equation must be used.
\end{itemize}

This is identical to the situation in Newtonian gravity: the Poisson equation
$\nabla^2\Phi = 4\pi G\rho$ is the quasi-static limit of the wave equation
$\Box\Phi = -4\pi G\rho$ (linearized GR in the weak-field limit). Nobody
questions using Poisson's equation for galaxies.


%=============================================================================
\section{Application to Quantum Experiments}
\label{sec:quantum}
%=============================================================================

\subsection{The key question}

The i32--i34 QG analyses used the constraint interpretation of $\psi$ to
argue that $\psi$ tracks each quantum branch independently (branch-dependent
classical gravity). For this to be self-consistent, the quasi-static limit
must be valid on the timescales and length scales of quantum experiments.

The critical parameter is:
\begin{equation}
\mathcal{R} \equiv \frac{t_\psi}{T_{\rm quantum}}
= \frac{d/c_s}{T_{\rm quantum}}
\end{equation}
where $d$ is the spatial separation of the superposed masses and
$T_{\rm quantum}$ is the relevant quantum timescale (measurement time,
coherence time, etc.).

If $\mathcal{R} \ll 1$, the $\psi$-field has ample time to ``track'' the
quantum state on each branch, and the constraint interpretation is valid.

\subsection{Page-Geilker experiment (1981)}

\begin{tcolorbox}[colback=blue!5!white, colframe=blue!60!black,
title=Page-Geilker: Quasi-Static Assessment]
\textbf{Setup}: A radioactive decay triggers a mass to move to position $L$
or $R$, separated by $d \sim 1$~cm $= 10^{-2}$~m.

\textbf{$\psi$-propagation time}:
\begin{equation}
t_\psi = \frac{d}{c_s} \sim \frac{10^{-2}\text{ m}}{c/\sqrt{3}}
= \frac{10^{-2}}{1.73 \times 10^8} \approx 5.8 \times 10^{-11}\text{ s}
\end{equation}

(Using $c_s \sim c/\sqrt{3}$ for a weak-field laboratory environment where
$\bar\chi \ll 1$.)

\textbf{Quantum/measurement timescale}: The Cavendish balance used by
Page \& Geilker responds on timescales $T_{\rm meas} \sim 10^{-2}$--$1$~s
(the torsion balance period). Even the fastest conceivable quantum measurement
gives $T_{\rm quantum} \gtrsim 10^{-8}$~s.

\textbf{Quasi-static ratio}:
\begin{equation}
\mathcal{R} = \frac{t_\psi}{T_{\rm quantum}}
\sim \frac{6 \times 10^{-11}}{10^{-8}} = 6 \times 10^{-3}
\end{equation}

\textbf{Conclusion}: The quasi-static limit is valid by a factor of $\sim
160$. The $\psi$-field has ample time to readjust to whichever quantum branch
is realized. The constraint interpretation is fully justified.
\end{tcolorbox}

\subsection{BMV experiment (Bose-Marletto-Vedral, proposed)}

\begin{tcolorbox}[colback=green!5!white, colframe=green!60!black,
title=BMV: Quasi-Static Assessment]
\textbf{Setup}: Two masses in spatial superposition, separated by
$d \sim 200~\mu$m $= 2 \times 10^{-4}$~m. The experiment seeks to detect
gravitationally-induced entanglement over a coherence time
$T_{\rm coh} \sim 1$--$2$~s.

\textbf{$\psi$-propagation time}:
\begin{equation}
t_\psi = \frac{d}{c_s} \sim \frac{2 \times 10^{-4}}{1.73 \times 10^8}
\approx 1.2 \times 10^{-12}\text{ s}
\end{equation}

\textbf{Quasi-static ratio}:
\begin{equation}
\mathcal{R} = \frac{t_\psi}{T_{\rm coh}}
\sim \frac{1.2 \times 10^{-12}}{1} = 1.2 \times 10^{-12}
\end{equation}

\textbf{Conclusion}: The quasi-static limit is valid by \textbf{twelve
orders of magnitude}. For all practical purposes, $\psi$ is an instantaneous
constraint in BMV.
\end{tcolorbox}

\subsection{Microscopic quantum superpositions (atoms, molecules)}

For a superposition of an atom at two locations separated by
$d \sim 1~\mu$m $= 10^{-6}$~m, with coherence time $T_{\rm coh} \sim 1$~s:
\begin{equation}
t_\psi \sim \frac{10^{-6}}{1.73 \times 10^8} \approx 5.8 \times 10^{-15}\text{ s}
\end{equation}
\begin{equation}
\mathcal{R} \sim 6 \times 10^{-15}
\end{equation}

Valid by fifteen orders of magnitude.

For a superposition on a \emph{nuclear} scale $d \sim 1$~fm $= 10^{-15}$~m:
\begin{equation}
t_\psi \sim 3 \times 10^{-24}\text{ s}
\end{equation}

This is comparable to the strong-interaction timescale. Even for
$T_{\rm quantum} \sim 10^{-23}$~s (nuclear transitions), $\mathcal{R} \sim 0.3$.
The quasi-static limit begins to become marginal only at nuclear/subnuclear
length scales with nuclear-timescale dynamics.

\subsection{Summary table}

\begin{center}
\renewcommand{\arraystretch}{1.3}
\begin{tabular}{lcccc}
\toprule
\textbf{Experiment} & $d$ & $T_{\rm quantum}$ & $t_\psi$ & $\mathcal{R}$ \\
\midrule
BMV (proposed) & $200~\mu$m & $\sim 1$~s & $\sim 10^{-12}$~s & $10^{-12}$ \\
Page-Geilker & $\sim 1$~cm & $\gtrsim 10^{-8}$~s & $\sim 6\times 10^{-11}$~s & $6\times 10^{-3}$ \\
Atom interferometry & $\sim 1~\mu$m & $\sim 1$~s & $\sim 6\times 10^{-15}$~s & $6\times 10^{-15}$ \\
Molecule interferometry & $\sim 100$~nm & $\sim 10^{-3}$~s & $\sim 6\times 10^{-16}$~s & $6\times 10^{-13}$ \\
Nuclear transition & $\sim 1$~fm & $\sim 10^{-23}$~s & $\sim 3\times 10^{-24}$~s & $\sim 0.3$ \\
\midrule
Galaxy rotation & $30$~kpc & $10^{8}$~yr & $10^{5}$~yr & $10^{-3}$ \\
Binary pulsar (orbital) & $\sim 10^9$~m & $10^4$~s & $\sim 6$~s & $6\times 10^{-4}$ \\
NS merger (final) & $\sim 30$~km & $\sim 10^{-3}$~s & $\sim 2\times 10^{-4}$~s & $\sim 0.2$ \\
\bottomrule
\end{tabular}
\end{center}


%=============================================================================
\section{Precise Comparison: $\psi$ vs.\ $A_0$}
\label{sec:comparison}
%=============================================================================

\subsection{Where the analogy holds exactly}

\begin{center}
\renewcommand{\arraystretch}{1.3}
\begin{tabular}{>{\raggedright}p{5cm}cc}
\toprule
\textbf{Property} & \textbf{$A_0$ in QED} & \textbf{$\psi$ in DFD} \\
\midrule
Free propagator ($X=0$) & Absent & Absent \\
Background propagator ($X>0$) & N/A (always linear) & Exists, with $c_s^2 = (1+\chi)/(3+\chi)$ \\
Determined by sources (static) & Gauss's law (linear) & AQUAL equation (nonlinear) \\
$\dot{A}_0^2$ / $\dot\psi^2$ in Lagrangian & Absent identically & Present via $P(X)$ kinetic term \\
Elliptic in quasi-static limit & Yes & Yes \\
Full equation hyperbolic & No (always elliptic) & Yes \\
Constraint on each time slice & Exact & Approximate (to accuracy $\mathcal{R}^2$) \\
\bottomrule
\end{tabular}
\end{center}

\subsection{The critical difference}

In QED (Coulomb gauge), $A_0$ is exactly non-dynamical. The Lagrangian
$\mathcal{L}_{\rm QED}$ contains $\dot{\mathbf{A}}^2$ but \textbf{not}
$\dot{A}_0^2$. Varying with respect to $A_0$ gives Gauss's law, which is
a constraint at every instant. This is not an approximation---it is an exact
property of the gauge-fixed theory.

In DFD, $\psi$ enters the Lagrangian through the k-essence kinetic function:
\begin{equation}
\mathcal{L}_{\rm DFD} \supset P(X) = P\!\left(-\frac{1}{2}\partial_\mu\psi\,\partial^\mu\psi\right)
\end{equation}
where $X = -\frac{1}{2}[(\dot\psi/c)^2 - |\nabla\psi|^2]$. The time
derivative $\dot\psi$ enters through $X$, giving a $\dot\psi^2$ contribution
to the equations of motion. Specifically, expanding $P(X)$ around the
spatial-gradient background:
\begin{equation}
P(X) = P(X_0) + P'(X_0)\delta X + \frac{1}{2}P''(X_0)(\delta X)^2 + \ldots
\end{equation}
where $X_0 = \frac{1}{2}|\nabla\bar\psi|^2$ and
$\delta X \ni -\frac{1}{2}(\dot\psi/c)^2$, so $\dot\psi^2$ terms are
present and generate genuine dynamics.

\subsection{Why the analogy is nevertheless correct in practice}

\begin{proposition}[Effective Constraint Status]
\label{prop:effective-constraint}
In DFD, $\psi$ is an \emph{effective} constraint field to accuracy
$\order{\mathcal{R}^2}$, where $\mathcal{R} = L/(c_s T)$ is the
quasi-static ratio. For all laboratory quantum experiments,
$\mathcal{R}^2 < 10^{-5}$. The corrections to the constraint
interpretation are negligible.
\end{proposition}

\begin{proof}
Write $\psi = \psi_{\rm QS}[\rho(t)] + \delta\psi$, where
$\psi_{\rm QS}[\rho(t)]$ is the quasi-static solution (satisfying the
elliptic equation~(\ref{eq:elliptic}) at each instant with the
instantaneous $\rho(t)$). Substituting into the full
equation~(\ref{eq:full-dfd}):
\begin{equation}
\frac{1}{c^2}\ddot\psi_{\rm QS}
+ \frac{1}{c^2}\ddot{\delta\psi}
- \nabla\cdot[\mu\nabla(\psi_{\rm QS}+\delta\psi)]
= -\frac{8\pi G}{c^2}\rho
\end{equation}
Since $\psi_{\rm QS}$ satisfies the elliptic equation, the spatial
derivatives cancel with the source, leaving:
\begin{equation}
\frac{1}{c^2}\ddot\psi_{\rm QS}
+ \frac{1}{c^2}\ddot{\delta\psi}
- \nabla\cdot[\mu\nabla\delta\psi]
- \text{(nonlinear cross-terms)} = 0
\end{equation}
The leading correction is $\delta\psi \sim (L^2/c_s^2 T^2)\psi_{\rm QS}
= \mathcal{R}^2\,\psi_{\rm QS}$.
\end{proof}


%=============================================================================
\section{When Does the Quasi-Static Limit Break Down?}
\label{sec:breakdown}
%=============================================================================

The quasi-static limit breaks down when $\mathcal{R} \equiv L/(c_s T)
\gtrsim 1$. This occurs in:

\subsection{Ultra-relativistic astrophysics}

\begin{enumerate}
\item \textbf{Gravitational wave emission}: For binary inspirals with
      $v/c \gtrsim 0.1$, the retardation of $\psi$ matters. The full
      hyperbolic equation must be solved. However, as established in i33,
      the tensor modes $h_{ij}^{TT}$ propagate at exactly $c$ on the flat
      background, and the $\psi$-perturbations are \emph{slaved} to the
      tensor/matter evolution (not independently radiated). The corrections
      to GW waveforms from $\psi$-retardation are suppressed by
      $(v/c)^2 \times (\mu_{\rm MOND~correction})$.

\item \textbf{Neutron star mergers}: During the final milliseconds,
      $\mathcal{R} \sim 0.2$. The quasi-static approximation receives
      $\sim 4\%$ corrections.

\item \textbf{Cosmological perturbation theory}: On horizon-scale modes
      ($L \sim c/H_0$), $\mathcal{R} \sim 1$ by definition. The full
      time-dependent equation must be used. This is the regime relevant
      for CMB physics (large-angle modes) and structure formation on the
      largest scales.
\end{enumerate}

\subsection{Cosmological context}

In FRW cosmology, the background $\nabla\bar\psi = 0$ (homogeneity),
so $\bar\chi = 0$ and the entire kinetic sector degenerates. This is
the cuscuton limit identified in i32. Perturbations around FRW must be
treated as perturbations of the full action, not of the quasi-static
equation. The cosmological perturbation theory of DFD requires special
treatment (see the DFD Cosmology paper for details).

\subsection{What does NOT break the quasi-static limit}

\begin{itemize}
\item \textbf{Quantum superpositions of microscopic masses}: Even for
      the fastest quantum processes (nuclear timescales $\sim 10^{-23}$~s),
      the $\psi$-propagation time across nuclear distances ($\sim 10^{-24}$~s)
      is still shorter. The quasi-static limit holds marginally. For all
      atomic and molecular superpositions, it holds by many orders of
      magnitude.

\item \textbf{Rapid quantum measurements}: A measurement that collapses
      a superposition ``instantaneously'' still has a physical timescale
      (the decoherence time). This is always $\gg t_\psi$ for any
      macroscopic or mesoscopic system.

\item \textbf{Quantum tunneling}: The tunneling time through a barrier
      of width $d$ is at minimum $d/c$ (by relativistic causality).
      The $\psi$-propagation time $d/c_s \leq \sqrt{3}\,d/c$ is at most
      $\sqrt{3}$ times longer. So even for tunneling, $\psi$ adjusts
      on a comparable timescale.
\end{itemize}


%=============================================================================
\section{The ``Slaving'' Argument}
\label{sec:slaving}
%=============================================================================

Because $c_s \leq c$, one might worry that $\psi$ perturbations ``lag behind''
matter. This worry is unfounded for the quasi-static regime, but worth
addressing precisely.

\subsection{Subluminal propagation and slaving}

In the Newtonian regime ($\bar\chi \gg 1$), $c_s \to c$. The $\psi$-field
propagates nearly as fast as light. Any change in the matter distribution
is communicated through $\psi$ at essentially the speed of light.

In the MOND regime ($\bar\chi \ll 1$), $c_s \to c/\sqrt{3}$. The
$\psi$-field propagates at $\sim 58\%$ of the speed of light. This is
slower than light, but still \emph{much faster than any matter motion}
in the non-relativistic regime.

The key point: for the $\psi$-field to ``track'' the matter distribution
(i.e., for the constraint interpretation to hold), we need $c_s \gg v_{\rm matter}$, not $c_s \geq c$. For all non-relativistic systems:
\begin{equation}
\frac{v_{\rm matter}}{c_s} \leq \sqrt{3}\,\frac{v_{\rm matter}}{c} \ll 1
\end{equation}

\subsection{Adiabatic following}

The situation is analogous to the Born-Oppenheimer approximation in molecular
physics: electrons ($\psi$) are ``fast'' degrees of freedom that adiabatically
follow the ``slow'' nuclear motion (matter). As long as the electronic
(field-propagation) timescale is much shorter than the nuclear (matter-evolution)
timescale, the electrons are effectively in their ground state for each nuclear
configuration. Corrections enter at order $(m_e/m_N)^{1/2} \sim \mathcal{R}$.

In DFD:
\begin{itemize}
\item ``Fast'' degree of freedom: $\psi$ with propagation speed $c_s \sim c$
\item ``Slow'' degree of freedom: matter with velocity $v \ll c$
\item Adiabaticity parameter: $\mathcal{R} = v/c_s \leq \sqrt{3}\,v/c$
\item The $\psi$-field adiabatically follows the matter distribution on
      each quantum branch.
\end{itemize}


%=============================================================================
\section{Corrected Characterization of $\psi$}
\label{sec:corrected}
%=============================================================================

\begin{tcolorbox}[colback=red!5!white, colframe=red!60!black,
title=Corrected Statement (Replaces the i32--i34 Characterization)]

The DFD $\psi$-field has a \textbf{dual character}:

\begin{enumerate}
\item \textbf{Mathematically}: $\psi$ satisfies a nonlinear hyperbolic wave
equation with sound speed $c_s^2 = (1+\chi)/(3+\chi) \in [1/3, 1)$.
It is a dynamical, propagating field---\textbf{not} literally a constraint
like $A_0$ in QED.

\item \textbf{Physically, in the quasi-static regime} ($L \ll c_s T$):
$\psi$ is effectively a constraint field, instantaneously determined by
the matter distribution to accuracy $\order{(L/c_s T)^2}$. This regime
encompasses:
\begin{itemize}
\item All galactic and cluster dynamics (standard MOND phenomenology)
\item All solar system tests
\item \textbf{All laboratory quantum gravity experiments} (Page-Geilker,
BMV, atom interferometry, etc.)
\item All conceivable quantum experiments with non-relativistic masses
\end{itemize}

\item \textbf{The i32--i34 QG programme is correct in substance}: For the
quantum experiments analyzed (Page-Geilker, BMV, entanglement witnesses),
calling $\psi$ a ``constraint field'' is accurate to better than one
part in $10^5$. The branch-dependent quantization scheme
$|\Psi\rangle = \sum_n c_n |n\rangle|\psi_n\rangle$, where $\psi_n$
solves the elliptic equation with source $\rho_n$, is correct in this
regime.

\item \textbf{The Coulomb analogy is an approximation, not an identity}:
$A_0$ is exactly non-dynamical (no $\dot{A}_0$ in the Lagrangian).
$\psi$ is approximately non-dynamical (the $\dot\psi$ contribution is
suppressed by $\mathcal{R}^2$). The analogy is quantitatively excellent
for all relevant applications but fails at the formal level. This must
be stated clearly in publications.

\item \textbf{Failure of the quasi-static limit}: The full wave equation
must be used for gravitational wave physics ($v/c \gtrsim 0.1$),
cosmological perturbations on horizon scales, and (marginally) for
nuclear-timescale processes at nuclear length scales.
\end{enumerate}
\end{tcolorbox}


%=============================================================================
\section{Implications for the Branch-Dependent Quantization}
\label{sec:implications}
%=============================================================================

\subsection{Does the wave-equation nature of $\psi$ spoil the QG programme?}

No. The branch-dependent quantization of i32--i34 required that $\psi$ be
determined by the matter state on each branch. This is satisfied in two
equivalent ways:

\textbf{Exact (full equation)}: On each branch $|n\rangle$ with matter
density $\rho_n(\mathbf{x},t)$, $\psi_n(\mathbf{x},t)$ is determined
by the retarded solution of the full hyperbolic equation with source
$\rho_n$. The retardation introduces a delay of $\order{L/c_s}$, which
is negligible for laboratory experiments.

\textbf{Approximate (quasi-static)}: On each branch, $\psi_n$ is
determined by the instantaneous solution of the elliptic equation with
source $\rho_n(t)$. This is accurate to $\order{\mathcal{R}^2}$.

Both give the same physics for quantum experiments. The branch-dependent
structure $|\Psi\rangle = \sum_n c_n|n\rangle|\psi_n\rangle$ is valid
in either formulation, because in both cases $\psi_n$ is a \emph{functional}
of $\rho_n$---it has no independent degrees of freedom (no free
$\psi$-quanta propagating independently of matter).

\subsection{Why there are no free $\psi$-quanta}

This is the deepest point, and it survives the correction. The i32 result
that the free propagator vanishes at $\bar\chi = 0$ remains true. On a
non-trivial background, $\psi$ perturbations propagate, but they are
\emph{sourced perturbations}---always tied to matter fluctuations.

The absence of free $\psi$-radiation (i33, confirmed by zero scalar GW
memory) means that $\psi$ never develops an independent quantum state.
It is always slaved to matter, whether instantaneously (quasi-static)
or retardedly (full equation). This is what makes the branch-dependent
quantization consistent.

\subsection{Comparison with the graviton}

In GR, the graviton is a freely propagating, independently quantized degree
of freedom ($h_{ij}^{TT}$ satisfies $\Box h_{ij}^{TT} = 0$ in vacuum).
In DFD, the tensor mode $h_{ij}^{TT}$ on the flat background is the analog
of the graviton and does propagate freely. But $\psi$ is \textbf{not} a
graviton---it is the scalar sector, and it has no free excitations.

This is structurally identical to the Coulomb sector of QED: $A_0$ has no
photon excitations, but $A_i^T$ does. The DFD $\psi$ is the gravitational
$A_0$; the tensor modes are the gravitational $A_i^T$.


%=============================================================================
\section{Explicit Verification: Page-Geilker Revisited}
\label{sec:page-geilker}
%=============================================================================

The Page-Geilker experiment is the most conceptually important test
of semiclassical gravity. The argument proceeds:

\begin{enumerate}
\item A radioactive source determines whether a mass moves left or right.
\item Before measurement, the quantum state is
$|\Psi\rangle = \frac{1}{\sqrt{2}}(|L\rangle + |R\rangle)$.
\item Semiclassical gravity predicts the gravitational field responds to
$\langle\hat\rho\rangle = \frac{1}{2}(\rho_L + \rho_R)$---the average.
\item The experiment found the field responds to the \emph{actual} outcome
($\rho_L$ or $\rho_R$), not the average.
\end{enumerate}

In DFD's branch-dependent scheme:
\begin{equation}
|\Psi\rangle = \frac{1}{\sqrt{2}}\bigl(|L\rangle|\psi_L\rangle
+ |R\rangle|\psi_R\rangle\bigr)
\end{equation}
where $\psi_L$ and $\psi_R$ are the quasi-static solutions for the
left and right configurations respectively.

The quasi-static approximation is justified because:
\begin{itemize}
\item The mass moves at $v \sim 1$~m/s. The $\psi$-field adjusts at
$c_s \sim 10^8$~m/s. The ratio $v/c_s \sim 10^{-8}$.
\item The Cavendish balance responds on timescale $T \sim 0.1$~s.
The $\psi$-propagation time across the balance ($\sim 0.1$~m) is
$\sim 10^{-9}$~s.
\item $\mathcal{R} \sim 10^{-8}$ for the mass motion;
$\mathcal{R} \sim 10^{-8}$ for the balance response.
\end{itemize}

The correction to the constraint interpretation is
$\order{\mathcal{R}^2} \sim 10^{-16}$---completely negligible.


%=============================================================================
\section{Explicit Verification: BMV Experiment}
\label{sec:bmv}
%=============================================================================

For the BMV experiment:
\begin{itemize}
\item Two diamond microspheres (mass $\sim 10^{-14}$~kg) are placed in
spatial superpositions with arm separation $d \sim 200~\mu$m.
\item The gravitational interaction over coherence time $T_{\rm coh} \sim 2$~s
is predicted to entangle the two masses.
\item If $\psi$ is a classical constraint field (branch-dependent), DFD
predicts entanglement occurs (because each branch has a different
$\psi$ mediating a different gravitational interaction).
\end{itemize}

Quasi-static verification:
\begin{equation}
t_\psi = \frac{d}{c_s} = \frac{2\times 10^{-4}}{1.73\times 10^8}
= 1.16 \times 10^{-12}\text{ s}
\end{equation}
\begin{equation}
\mathcal{R} = \frac{t_\psi}{T_{\rm coh}} = 5.8 \times 10^{-13}
\end{equation}

The constraint approximation is valid to one part in $10^{24}$.

\begin{remark}
The BMV experiment cannot distinguish DFD's branch-dependent classical
gravity from full quantum gravity (with a propagating graviton), because
both predict entanglement. The experiment would rule out semiclassical
gravity ($\langle\hat{T}_{\mu\nu}\rangle$-based), which DFD already
excludes on theoretical grounds (i33, Theorem 3).
\end{remark}


%=============================================================================
\section{Answers to the Six Questions}
\label{sec:answers}
%=============================================================================

\begin{enumerate}
\item \textbf{When is the quasi-static limit valid?}

Whenever $L \ll c_s T$, i.e., whenever the system is non-relativistic
($v/c \ll 1$). This includes all galactic dynamics, all solar system
physics, and all laboratory experiments. The $\psi$-field then satisfies
the elliptic AQUAL equation and is instantaneously determined by the
matter distribution---it behaves as a constraint field.

\item \textbf{Does $\psi$ propagate subluminally? Is it ``slaved'' to
matter?}

Yes and yes. The sound speed satisfies $c/\sqrt{3} \leq c_s < c$
everywhere (on non-trivial backgrounds). Since $c_s \gg v_{\rm matter}$
for all non-relativistic systems, $\psi$ adiabatically follows the
matter distribution. It is ``slaved'' in the same sense that the
Coulomb field of a slowly-moving charge is slaved to the charge position.
The $\psi$-perturbations are always ``catching up'' to matter changes
on a timescale $t_\psi = L/c_s \ll T_{\rm matter}$.

\item \textbf{Is the quasi-static limit valid for quantum superpositions
at microscopic distances?}

Yes. For atomic separations ($\sim\mu$m), $t_\psi \sim 10^{-15}$~s,
while quantum coherence times are $\sim 10^{-3}$--$1$~s. The ratio
$\mathcal{R} \sim 10^{-12}$--$10^{-15}$. The quasi-static limit is
spectacularly valid.

\item \textbf{Is the Page-Geilker analysis correct?}

Yes. The $\psi$-propagation time across centimeters is
$\sim 6 \times 10^{-11}$~s. The measurement timescale is
$\gtrsim 10^{-8}$~s. The quasi-static ratio is $\mathcal{R} \sim 10^{-3}$,
and the correction is $\mathcal{R}^2 \sim 10^{-6}$. The constraint
interpretation is accurate to one part per million.

\item \textbf{Is the BMV analysis correct?}

Yes. The quasi-static ratio is $\mathcal{R} \sim 10^{-12}$, valid by
twelve orders of magnitude. The $\psi$-field is effectively an
instantaneous constraint for BMV.

\item \textbf{When does the quasi-static limit break down?}

For ultra-relativistic processes:
\begin{itemize}
\item Gravitational wave generation/propagation ($v/c \gtrsim 0.1$)
\item Cosmological perturbations on horizon scales ($L \sim c/H_0$)
\item Compact object mergers (final inspiral, $v \sim 0.3\,c$)
\item (Marginally) Nuclear processes at nuclear length scales
\end{itemize}

For \textbf{all} laboratory quantum experiments with non-relativistic
masses, $\psi$ is effectively a constraint field. This is the correct
characterization.
\end{enumerate}


%=============================================================================
\section{Conclusion: The Precise Statement}
\label{sec:conclusion}
%=============================================================================

The DFD $\psi$-field is \textbf{not} a constraint field in the strict sense
of $A_0$ in Coulomb-gauge QED (which has no time derivative in the
Lagrangian). It is a dynamical field governed by a nonlinear wave equation
with subluminal sound speed $c/\sqrt{3} \leq c_s < c$.

However, for all non-relativistic systems---which includes every laboratory
quantum experiment ever proposed or conceived---the quasi-static limit
renders $\psi$ an \emph{effective} constraint, determined by the
instantaneous matter distribution to extraordinary accuracy
($\mathcal{R}^2 < 10^{-5}$ for Page-Geilker; $\mathcal{R}^2 < 10^{-24}$
for BMV).

The i32--i34 QG programme's central results---branch-dependent classical
gravity, resolution of semiclassical paradoxes, compatibility with Page-Geilker,
prediction of gravitationally-induced entanglement---\textbf{all stand},
because they require only the quasi-static limit, which is justified by the
enormous hierarchy between $c_s$ and $v_{\rm matter}$ in all relevant
experiments.

The key refinement for publications: replace ``$\psi$ is a constraint field''
with ``$\psi$ is an effective constraint field in the quasi-static limit,
which is valid for all non-relativistic systems.'' This is more precise and
forestalls the obvious objection that the full equation is hyperbolic.

\begin{tcolorbox}[colback=green!5!white, colframe=dfdgreen,
title=Bottom Line]
The dual character of $\psi$ (wave equation + effective constraint) is a
\emph{strength}, not a weakness. It means DFD can describe both:
\begin{itemize}
\item Static/quasi-static gravity (galactic dynamics, quantum experiments)
via the elliptic constraint equation, and
\item Radiative gravity (gravitational waves, cosmological perturbations)
via the full hyperbolic equation.
\end{itemize}
This is richer than both pure-constraint theories (which cannot describe
GWs) and standard scalar-tensor theories (which have freely propagating
scalar modes that are observationally constrained). DFD's $\psi$ propagates
but is never free---always slaved to matter. This is the structural analog
of the Coulomb sector in QED, elevated to the full nonlinear theory.
\end{tcolorbox}


%=============================================================================
% References
%=============================================================================
\begin{thebibliography}{99}

\bibitem{Alcock2025DFD}
G.~T.~Alcock, ``Density Field Dynamics: A Complete Unified Theory,''
v10.8 (2025).

\bibitem{BekensteinMilgrom1984}
J.~Bekenstein and M.~Milgrom,
``Does the missing mass problem signal the breakdown of Newtonian gravity?''
\textit{Astrophys.~J.}~\textbf{286}, 7 (1984).

\bibitem{PageGeilker1981}
D.~N.~Page and C.~D.~Geilker,
``Indirect evidence for quantum gravity,''
\textit{Phys.~Rev.~Lett.}~\textbf{47}, 979 (1981).

\bibitem{BMV2017}
S.~Bose \textit{et al.},
``Spin entanglement witness for quantum gravity,''
\textit{Phys.~Rev.~Lett.}~\textbf{119}, 240401 (2017).

\bibitem{Marletto2017}
C.~Marletto and V.~Vedral,
``Gravitationally induced entanglement between two massive particles is
sufficient evidence of quantum effects of gravity,''
\textit{Phys.~Rev.~Lett.}~\textbf{119}, 240402 (2017).

\bibitem{Afshordi2006}
N.~Afshordi, D.~J.~H.~Chung, and G.~Geshnizjani,
``Cuscuton: A causal field theory with an infinite speed of sound,''
\textit{Phys.~Rev.~D}~\textbf{75}, 083513 (2007).

\bibitem{Fragkos2025}
V.~Fragkos \textit{et al.},
``Classical theories of gravity produce entanglement,''
\textit{Nature}~\textbf{639}, 531 (2025).

\bibitem{BabichevEsposito2025}
E.~Babichev and G.~Esposito-Far\`{e}se,
``Retarded solutions of nonlinear wave equations in curved spacetime,''
arXiv:2409.16108 (2024).

\end{thebibliography}

\end{document}
